% ===================================================================
%  పట్టిక సంఖ్య 11 — నగరం, దాని కట్టడాలు. అగ్ని ప్రమాదం.
%
%  Traduction, faite en 2026, en telougou d'aujourd'hui, sur l'ido de
%  texto/io. Regles de la colonne : voir 10-pattika-01.tex. Deux
%  parties, deux numerotations : les cles portent c1 et c2. Le bloc
%  « apar » porte la mention de serie, le filet, le numero et le
%  titre, comme en hindi et en marathe.
%
%  LES INSTITUTIONS SE DECRIVENT, ELLES NE SE TRANSPOSENT PAS. C'est
%  le point delicat de ce tableau : une prefecture, une Bourse, une
%  ecole normale de 1926 n'ont pas d'equivalent administratif en
%  Andhra, et leur donner le nom de l'institution indienne la plus
%  proche — కలెక్టరు కార్యాలయం pour la prefecture — aurait deplace
%  la CHOSE, ce que la regle interdit. La colonne ecrit donc la
%  fonction : పరిపాలన భవనం l'edifice de l'administration (15),
%  వర్తక భవనం l'edifice du negoce (14), ఉపాధ్యాయ శిక్షణ కళాశాల
%  l'ecole ou l'on forme les maitres (41). C'est exactement ce
%  qu'avaient fait le hindi (प्रशासन-भवन, व्यापार-भवन) et le
%  marathe avant lui : trois langues, une seule solution, et ce
%  n'est pas une coincidence — c'est la seule qui ne mente pas.
%
%  LE MOT-OUTIL NE SE GRASSE PAS, MEME QUAND LE FAC-SIMILE LE GRASSE.
%  Le bloc c1-09-1 de l'ido met en gras « quale », « La » et « ol » au
%  milieu de groupes substantifs. Les colonnes traduites degrassent
%  les mots-outils depuis le tableau 1 ; celle-ci garde le gras sur
%  les groupes — ట్రాముల లాగుడు వ్యవస్థ, గుర్రాలు లాగే, నగరం ఒక
%  చివరనుంచి మరో చివరకు, చౌకగా — et le laisse tomber sur les
%  articles et les conjonctions.
%
%  LE (81) EST DONNE DEUX FOIS PAR LE FAC-SIMILE, aux musiciens et au
%  violon du meme bloc c2-06-1. C'est une faute de l'edition ido, et
%  LA SOURCE NE BOUGE PAS : la traduction porte les deux, comme les
%  trente-cinq colonnes qui precedent.
%
%  ONZE RETOURNEMENTS, DIX SUR LES VINGT PAIRES DE parigi.py et UN
%  QU'ELLE N'A PAS VU. Les dix vus sont des relatives et des genitifs
%  ordinaires — « pordo (8) dil kastelo (7) », « stradeto (11) qua
%  duktas al komon-domo (12) », « kupolo (40) di la hospitalo (39) »,
%  « muzikisti (81) di l'orkestro (80) ». Les dix non retournees se
%  rangent dans les trois cases du tableau 9 : enumeration, renvoi de
%  gauche deja modificateur (« gasiferio (5) di qua la kameno (6) » —
%  le telougou dit « la cheminee de l'usine a gaz » dans cet ordre),
%  frontiere de phrase.
%
%  ET LE ONZIEME EST UN DEUXIEME ECHAPPEMENT PAR PORTEE, deux
%  tableaux de suite. « Tribunalo (22) avan qua extensas su agreabla
%  gardeno publika (26) » : le relatif y est, mais 47 caracteres
%  separent les deux renvois et PORTEO vaut 45. Au tableau 10 le meme
%  echappement mesurait 54. Deux points, donc, et la tentation
%  d'elargir. On a mesure avant de ceder : porter la portee a 55
%  ferait passer le livret de 386 a 465 paires — vingt pour cent de
%  lignes en plus pour rattraper deux inversions que renvoji.py
%  attrape de toute facon. La mesure est reportee dans l'entete de
%  parigi.py, avec les deux distances ; le seuil ne bouge pas.
% ===================================================================

%%K t11-apar-1 apar
\VUsaut{2.0mm}
\VUcentre{13.2pt}{90}{మూడవ శ్రేణి}
\VUsaut{3.0mm}
\VUfilet{16mm}
\VUsaut{5.6mm}
\VUtitre{15.0pt}{60}{పట్టిక సంఖ్య 11}
\VUsaut{1.4mm}
\VUpk{11.4pt}{40}{నగరం, దాని కట్టడాలు. --- అగ్ని ప్రమాదం.}
\VUsaut{2.2mm}

%%K t11-tit-1 sub
\VUpk{10.2pt}{40}{శివారు పేట.}
\VUsaut{3.0mm}

%%K t11-c1-01-1 p
1. --- మహాసముద్రానికి వంద కిలోమీటర్ల దూరంలో ఉన్న పెద్ద నగరంలో నేను ఉంటాను; అయినా అది
మొదటి తరగతి \VUgras{ఓడరేవు} \textsuperscript{(1)}. దాని పక్కనే పారే నది నాలుగు వందల
మీటర్ల వెడల్పు, అది చాలా అందమైన లంగరు మడుగును ఏర్పరుస్తుంది.

%%K t11-c1-02-1 p
2. --- \VUgras{రేవు గట్ల} \textsuperscript{(2)} వెంబడి ఉన్న నగర భాగమంతా పూర్తిగా సముద్ర,
పారిశ్రామిక రూపం కలిగి ఉంది; అది ఒక \VUgras{శివారు పేటగా} \textsuperscript{(3)}
ఏర్పడింది — అక్కడ నావికులూ కూలీలూ మాత్రమే ఉంటారు. వీళ్ళు \VUgras{కర్మాగారాల్లో}
\textsuperscript{(4)}, \VUgras{గ్యాస్ కర్మాగారంలో} \textsuperscript{(5)} పనిచేస్తారు; ఆ
కర్మాగారపు ఎత్తైన \VUgras{పొగగొట్టం} \textsuperscript{(6)}, \VUgras{గ్యాస్ గోళం} చాలా
దూరంనుంచే కనిపిస్తాయి.

%%K t11-c1-03-1 p
3. --- మధ్యయుగంలో నగరానికి కోటబందీ ఉండేది; ఇప్పటికీ దాని ప్రాకారాల ఆనవాళ్ళు కొన్ని
దొరుకుతాయి — ముఖ్యంగా \VUgras{ద్వారం} \textsuperscript{(8)}, అది పాత
\VUgras{పటిష్ట కోటది} \textsuperscript{(7)}; దాని రెండు \VUgras{బురుజులు}
\textsuperscript{(9)} ఇరువైపులా నిలబడి ఉన్నాయి, ఇప్పుడు అవి \VUgras{చెరసాలగా} వాడుకలో
ఉన్నాయి.

%%K t11-c1-04-1 p
4. --- చాలా పాతదిగా కనిపించే ఈ \VUgras{పేటలో} ఒకే ఒక కట్టడం సాపేక్షంగా కొత్తది : అదే
\VUgras{సుంకపు కార్యాలయం} \textsuperscript{(10)}. అది రేవు గట్టుకూ \VUgras{సందుకూ}
\textsuperscript{(11)} మధ్య ఉంది; ఆ సందు పాత \VUgras{పురపాలక భవనం} \textsuperscript{(12)}
దాకా వెళ్తుంది. ఆ చివరి కట్టడాన్ని సులభంగా గుర్తుపట్టవచ్చు — దాని భారీ
\VUgras{గంట గోపురం} \textsuperscript{(13)} పాత నగరమంతటిపై పెత్తనం చేస్తుంది.

%%K t11-c1-05-1 p
5. --- వీధికి అవతలి వైపున సుంకపు కార్యాలయం శైలిలోనే కట్టిన భవనం ఒకటి \VUgras{నిలబడి}
ఉంది : అదే \VUgras{వర్తక భవనం} \textsuperscript{(14)}. దాని ఒక ముఖం చిన్న చౌక వైపు
చూస్తుంది, అక్కడే \VUgras{పరిపాలన భవనం} \textsuperscript{(15)} ఉంది. నిజానికి మా నగరం
మహారాజధాని కాకపోయినా, ముఖ్యమైన జిల్లా కేంద్రం.

%%K t11-tit-2 sub
\VUsaut{5.0mm}
\VUpk{10.2pt}{40}{నగరపు పై భాగం.}
\VUsaut{3.4mm}

%%K t11-c1-06-1 p
6. --- చాలా పెద్దదైన మా దండు, విశాలమైన, ఆరోగ్యకరమైన \VUgras{సైనిక శిబిరాల్లో}
\textsuperscript{(16)} ఉంటుంది; అవి \VUgras{నగర ఉద్యానవనం} \textsuperscript{(17)} ఇనుప
కంచె దాకా విస్తరించి ఉన్నాయి. ఆ ఉద్యానవనానికి వెళ్ళే \VUgras{వెడల్పు రోడ్డు}
\textsuperscript{(19)} \VUgras{పొదుపు బ్యాంకు} \textsuperscript{(18)} వెనుకనుంచి వెళ్ళి
\VUgras{ప్రధాన చర్చి} \textsuperscript{(20)} చౌక దగ్గర ముగుస్తుంది.

%%K t11-c1-07-1 p
7. --- ఆ కట్టడాలన్నీ ఒక చిన్న కొండపై అంచెలంచెలుగా విస్తరించి ఉన్నాయి; అక్కడే మా విశాలమైన
\VUgras{ఉన్నత పాఠశాల} \textsuperscript{(21)} కూడా ఉంది.

%%K t11-c1-07-2 p
కొంచెం వాలుగా ఉన్న దారి వెంట దిగితే — అది పెద్ద వీధిలో కలుస్తుంది — \VUgras{న్యాయస్థానం}
\textsuperscript{(22)} దగ్గరికి వస్తాం; దాని ముందు ఆహ్లాదకరమైన \VUgras{ప్రజా ఉద్యానం}
\textsuperscript{(26)} పరుచుకుని ఉంది. ఈ \VUgras{ఉద్యాన చౌక} పెద్దదేమీ కాదు, కానీ నగరం
నడిబొడ్డున పచ్చదనపు, చల్లదనపు ఒయాసిస్సును చేస్తుంది. అందుకే నగరవాసులు దానికి బాగా
వస్తుంటారు — \VUgras{కాఫీ దుకాణం} \textsuperscript{(25)} గుడారం కింద కూర్చునేందుకు,
గురువారం, ఆదివారం సైనిక సంగీతకారులు వాయించేది వినేందుకు; వాళ్ళు పచ్చిక బయళ్ళకూ పూల
గుత్తులకూ మధ్య పెట్టిన \VUgras{సంగీత మంటపంలో} \textsuperscript{(27)} వాయిస్తారు.

%%K t11-c1-08-1 p
8. --- ఆ పచ్చిక బయళ్ళలో ఒకదానిపై కంచు \VUgras{విగ్రహం} \textsuperscript{(28)} నిలబడి ఉంది
— మా నగరవాసుల్లో ఒకరి జ్ఞాపకార్థం కట్టిన స్మారకం అది; ఆయన తన జన్మనగరానికి గొప్ప సేవలు
చేశారు.

%%K t11-tit-3 sub
\VUsaut{5.0mm}
\VUpk{10.2pt}{40}{రవాణా సాధనాలు.}
\VUsaut{3.4mm}

%%K t11-c1-09-1 p
9. --- ఇప్పటిదాకా మా నగరానికి విద్యుత్ వెలుతురు లేదు, \VUgras{గ్యాస్ దీపాలు}
\textsuperscript{(29)} మాత్రమే ఉన్నాయి. కానీ \VUgras{ఇక్కడి పత్రికలు} ఈ వెలుతురు
వ్యవస్థను మెరుగుపరచాలని ప్రచారం చేస్తున్నాయి — \VUgras{ట్రాముల లాగుడు వ్యవస్థను}
ఇప్పటికే మెరుగుపరచినట్టుగానే. \VUgras{గుర్రాలు లాగే} ఆ పాత బళ్ళ స్థానంలో ఆచరణలో
\VUgras{విద్యుత్ ట్రాములు} \textsuperscript{(31)} వచ్చేశాయి; అవి ప్రయాణికులను
\VUgras{నగరం ఒక చివరనుంచి మరో చివరకు} చాలా \VUgras{చౌకగా} వేగంగా చేరుస్తాయి.

%%K t11-c1-10-1 p
10. --- న్యాయస్థానం దగ్గరే \VUgras{బ్యాంకు} \textsuperscript{(32)} ఉంది, అక్కడ వాణిజ్య
పత్రాలను మార్చుకుంటారు. \VUgras{బ్యాంకు వసూలుదారులు} \textsuperscript{(33)} ఇళ్ళకే వెళ్ళి
బాకీలు వసూలు చేస్తారు.

%%K t11-tit-4 sub
\VUsaut{5.0mm}
\VUpk{10.2pt}{40}{కళ, విద్య.}
\VUsaut{3.4mm}

%%K t11-c1-11-1 p
11. --- దగ్గరలో నిలబడిన అందమైన కట్టడం \VUgras{వస్తు ప్రదర్శనశాల}. అందులోనే నగర
\VUgras{గ్రంథాలయం} \textsuperscript{(34)}, అందమైన \VUgras{ప్రకృతి శాస్త్ర సేకరణలు}
\textsuperscript{(35)} (ప్రకృతి ప్రదర్శనశాల), ముద్దైన \VUgras{చిత్రశాల}
\textsuperscript{(36)} ఉన్నాయి; అది శాశ్వతమైన \VUgras{ప్రదర్శన స్థలం} అవుతుంది.

%%K t11-c1-11-2 p
వారానికి రెండుసార్లు అది ప్రజలకు తెరుస్తారు; కానీ బయటివాళ్ళు \VUgras{కావలివాడికి}
\textsuperscript{(37)} చిన్న మొత్తం ఇచ్చి ఏ రోజైనా చూడవచ్చు. \VUgras{చిత్రకారులు}
\textsuperscript{(38)} అక్కడికి పని చేసుకునేందుకు వస్తారు. తమ చట్రాల \VUgras{ముందు},
చేతిలో \VUgras{రంగుల పలకతో}, ప్రసిద్ధ చిత్రాలను నకలు తీస్తూ వాళ్ళు కనిపిస్తారు.

%%K t11-c1-12-1 p
12. --- ప్రదర్శనశాల వెనుక ఎత్తు మీద \VUgras{గోపురం} \textsuperscript{(40)} కనిపిస్తుంది
— \VUgras{ఆసుపత్రిది} \textsuperscript{(39)} అది —, ఇంకా
\VUgras{ఉపాధ్యాయ శిక్షణ కళాశాల} \textsuperscript{(41)} కట్టడాలు కనిపిస్తాయి; మా పురపాలక
పాఠశాలల ఉపాధ్యాయులకు అక్కడే శిక్షణ ఇచ్చేవారు.

%%K t11-c1-12-1 p suite
నాటకశాల చౌకలో \VUgras{తపాలా కార్యాలయం} \textsuperscript{(43)} ఉంది —
\VUgras{ప్రైవేటు ఇంట్లో} \textsuperscript{(42)} ఏర్పాటుచేసినది.

%%K t11-tit-5 sub
\VUsaut{6.0mm}
\VUpk{11.4pt}{40}{అగ్ని ప్రమాదం.}
\VUsaut{3.0mm}

%%K t11-tit-6 sub
\VUpk{10.2pt}{40}{మంటలు! అనే కేక.}
\VUsaut{3.4mm}

%%K t11-c2-01-1 p
1. --- నగరమంతా భయంకరమైన వార్త పాకుతోంది : నాటకశాలలో ఇప్పుడే మంటలు లేచాయి! నేను
\VUgras{చౌకకు} \textsuperscript{(96)} చేరుకునేసరికే జనం \VUgras{పాదచారుల దారిపైనా}
\textsuperscript{(46)}, \VUgras{రోడ్డుపైనా} \textsuperscript{(47)} గుమిగూడారు.
\VUgras{తపాలా గుమాస్తాలు} \textsuperscript{(44)}, \VUgras{ఉత్తరాలు పంచేవాళ్ళు}
\textsuperscript{(45)} తపాలా కార్యాలయంలోంచి బయటికి వస్తున్నారు. ఒక
\VUgras{ట్రాము డ్రైవరు} \textsuperscript{(48)} తన బండిని ఆపాల్సి వచ్చింది — నగర
\VUgras{అంబులెన్సు బండిని} \textsuperscript{(50)} పోనివ్వడానికి; అది
\VUgras{శస్త్రవైద్యుడితో} \textsuperscript{(52)}, \VUgras{నర్సుతో} \textsuperscript{(51)}
వస్తోంది — \VUgras{గాయపడినవాడిని} \textsuperscript{(53)} చూసేందుకు; అతన్ని
\VUgras{స్ట్రెచరుపై} \textsuperscript{(54)} \VUgras{వార్తాపత్రిక దుకాణం}
\textsuperscript{(55)} దగ్గరికి తెచ్చారు.

%%K t11-c2-02-1 p
2. --- కవాతునుంచి తిరిగి వస్తున్న పదాతిదళ బృందం ఒకటి ఆగిపోయింది — గుర్రాలపై ఉన్న
\VUgras{నగర రక్షక భటులతో} \textsuperscript{(61)} కలిసి క్రమశిక్షణ కాపాడేందుకు.
\VUgras{గుర్రపు రౌతుల} గుర్రాలు రెండు కాళ్ళపై లేస్తున్నాయి; \VUgras{సైనికుల}
\textsuperscript{(57)} కంటే, \VUgras{జెండార్ముల} \textsuperscript{(63)} కంటే వాళ్ళు
\VUgras{ఉత్సుకులను} \textsuperscript{(56)} బాగా వెనక్కి నెడతారు. అయినా జెండార్ములకు పని
ఎక్కువే. వాళ్ళలో ఒకడు \VUgras{పోలీసు సిపాయిలకు} \textsuperscript{(64)} చెబుతున్నాడు —
ఇంకో \VUgras{గాయపడినవాడిని} \textsuperscript{(65)} ఎక్కడికి తీసుకెళ్ళాలో; అతను నిచ్చెన
మీదినుంచి పడిన ధైర్యవంతుడైన అగ్నిమాపక సిబ్బంది.

%%K t11-c2-03-1 p
3. --- \VUgras{అధికారి} \textsuperscript{(66)} కచ్చితంగా చెబుతున్నాడు — వస్తున్న
\VUgras{ఆవిరి పంపును} \textsuperscript{(67)} ఎక్కడ పెట్టాలో. ఆ శక్తిమంతమైన యంత్రం కోసం
ఎదురుచూస్తూ, ఆయన హడావిడిగా \VUgras{చేతి పంపును} \textsuperscript{(68)} పెట్టించాడు; దాని
పొడవాటి \VUgras{గొట్టం} \textsuperscript{(69)} మంటలపై వరదలా నీళ్ళు చిమ్ముతుంది.

ఆరుగురు అగ్నిమాపక సిబ్బంది దాన్ని నడుపుతున్నారు; ఈలోగా వాళ్ళ సహచరుడు \VUgras{ముక్కును}
\textsuperscript{(70)} పట్టుకుని \VUgras{నీటి ధారను} \textsuperscript{(71)} మంట నడిబొడ్డుకు
గురిపెడుతున్నాడు.

%%K t11-c2-04-1 p
4. --- నాటకశాలకు అవతలి వైపున పెద్ద \VUgras{నిచ్చెనను} \textsuperscript{(72)} ఆనించారు.
దాని వల్ల అగ్నిమాపక సిబ్బంది మంటల దగ్గరికి చేరుకోగలుగుతారు — ఆ మంటలు నల్లని
\VUgras{పొగ} \textsuperscript{(75)} సుడుల మధ్యనుంచి ఎగిసిపడుతున్నాయి.

%%K t11-tit-7 sub
\VUsaut{5.0mm}
\VUpk{10.2pt}{40}{ధైర్య కృత్యాలు.}
\VUsaut{3.4mm}

%%K t11-c2-05-1 p
5. --- \VUgras{మంట} \textsuperscript{(74)} \VUgras{నాటకశాల} \textsuperscript{(73)}
విశ్రాంతి గదిలో పుట్టింది — \VUgras{సంగీత నాటకపు} సాధన జరుగుతున్నప్పుడు; దాని మొదటి
\VUgras{ప్రదర్శన} రేపు జరగాల్సి ఉండేది. అదృష్టవశాత్తూ ప్రేక్షక మందిరంలో కొద్దిమంది
\VUgras{ప్రేక్షకులే} \textsuperscript{(76)} ఉన్నారు. పాపం వాళ్ళు \VUgras{బాల్కనీలో}
\textsuperscript{(77)} ఆశ్రయం పొందారు; అక్కడికే ధైర్యవంతులైన
\VUgras{అగ్నిమాపక సిబ్బంది} \textsuperscript{(86)} నిచ్చెనతో, తాళ్ళతో సాయానికి వస్తున్నారు.

%%K t11-c2-06-1 p
6. --- \VUgras{స్తంభ వసారా} \textsuperscript{(78)} కింద — అక్కడే
\VUgras{ప్రదర్శన ప్రకటన} \textsuperscript{(79)} ఆ రాత్రి ఆటను తెలియజేస్తుంది —
\VUgras{సంగీతకారులు} \textsuperscript{(81)} పారిపోవడం కనిపిస్తుంది; \VUgras{వాద్యబృందం}
\textsuperscript{(80)} వాళ్ళే, తమ వాద్యాలు మోసుకుంటూ : \VUgras{వయొలిన్}
\textsuperscript{(81)}, \VUgras{వేణువు} \textsuperscript{(82)}, \VUgras{చెలో}
\textsuperscript{(83)}, \VUgras{ట్రాంబోన్} \textsuperscript{(84)}, \VUgras{డోలు}
\textsuperscript{(85)} ఇత్యాదులు. \VUgras{సామగ్రిలో} \textsuperscript{(87)} కొంత భాగాన్ని
కాపాడేందుకు ప్రయత్నిస్తున్నారు.

%%K t11-c2-07-1 p
7. --- \VUgras{నటులూ} \textsuperscript{(89)} \VUgras{నటీమణులూ} \textsuperscript{(88)},
\VUgras{గాయకులూ} \VUgras{గాయనీమణులూ} తమ రంగస్థల \VUgras{దుస్తుల్లోనే}
\textsuperscript{(90)} పారిపోవాల్సి వచ్చింది. ప్రధాన గాయకుడు ఒక
\VUgras{పత్రికా విలేకరికి} \textsuperscript{(92)} ఇదంతా ఎలా జరిగిందో వివరిస్తున్నాడు.
\VUgras{ఆ విలేకరి} మొదటి క్షణంనుంచే అధికారులతో పాటు పరుగెత్తుకొచ్చాడు :
\VUgras{పరిపాలనాధికారి} \textsuperscript{(91)}, \VUgras{నగరాధ్యక్షుడు}
\textsuperscript{(93)}, \VUgras{పోలీసు కమిషనరు} \textsuperscript{(94)},
\VUgras{న్యాయాధికారి} \textsuperscript{(95)} — వీళ్ళు బాధ్యత ఎవరిదో తేల్చేందుకు విచారణ
చేస్తారు.

\VUsaut{6.0mm}
\VUfilet{22mm}
